%% PI-AI-MM manuscript -- first full draft, 2026-11-09.
%% Template: Springer Nature sn-jnl (option sn-nature), as shipped in this folder (sn-article.tex is the untouched template example).
%% Tables are \input from ../tables and figures from ../figures (built by ../scripts/make_tables.py / make_figures.py) -- never retype numbers here
%% that live in a table. Every number in the text is traceable to ../results/ (see ../README.md, "Where each claim comes from").
%% Open items for the authors are marked \todo{...} and listed in DRAFT_NOTES.md.

\documentclass[pdflatex,sn-nature]{sn-jnl}

\usepackage{graphicx}%
\usepackage{multirow}%
\usepackage{amsmath,amssymb,amsfonts}%
\usepackage{amsthm}%
\usepackage{mathrsfs}%
\usepackage[title]{appendix}%
\usepackage{xcolor}%
\usepackage{textcomp}%
\usepackage{manyfoot}%
\usepackage{booktabs}%
\usepackage{url}%

\newcommand{\todo}[1]{\textcolor{red}{[TODO: #1]}}
\newcommand{\muB}{\ensuremath{\mu_\mathrm{B}}}
\newcommand{\Tc}{\ensuremath{T_\mathrm{C}}}

\raggedbottom

\begin{document}

\title[Accuracy is not reliability in magnetic-materials ML]{Accuracy is not reliability: physics descriptors, calibration and chemical generalization in machine learning of magnetic materials}

\author*[1]{\fnm{Mujahidul} \sur{Islam}}\email{mujahidmahfuz@gmail.com}

\author[2]{\fnm{Gayan P.} \sur{Hettiarachchi}}

\affil*[1]{\orgdiv{\todo{department}}, \orgname{\todo{institution}}, \orgaddress{\city{\todo{city}}, \country{\todo{country}}}}

\affil[2]{\orgname{Tokyo International University}, \orgaddress{\city{Kawagoe}, \state{Saitama}, \country{Japan}} \todo{confirm department and address}}

\abstract{Machine-learning models of magnetic materials are usually judged by test-set accuracy, and that accuracy is often taken to stand for a model's general quality. We tested that assumption on a curated, multi-source dataset of inorganic magnetic materials (magnetic ordering, 79,125 records; Curie temperature, 29,740; DFT magnetic moment per cell, 87,513). We asked three questions. Do physics descriptors from density functional theory improve prediction? Are the uncertainties calibrated? Do the predictions generalize to unfamiliar chemistry? Composition-only gradient-boosted models were accurate (ordering macro-F1 0.73; Curie temperature mean absolute error 76 K, $R^2$ 0.78). Adding DFT descriptors on identical materials and splits gave no practical gain for ordering or Curie temperature. For the moment, the small per-cell gain was fully explained by unit-cell size. Conformal intervals reached 90\% marginal coverage only with difficulty-normalized scores. Marginal ordering sets covered antiferromagnets and ferrimagnets only 57\% and 55\% of the time, against about 90\% for class-conditional sets. On an audited external benchmark of experimental Curie temperatures, compositions the model had never seen had an error of 103 K [79, 154]. In chemical systems absent from training, $R^2$ fell to $-0.17$ [$-0.93$, 0.66], and error grew with compositional distance. Accuracy, gains from physics descriptors, calibration and generalization are therefore separate properties, and each has to be measured before it is claimed.}

\keywords{magnetic materials, Curie temperature, magnetic ordering, machine learning, conformal prediction, uncertainty quantification, out-of-distribution generalization, reproducibility}

\maketitle

%% ===================================================================================================================
\section{Introduction}\label{sec:intro}

Magnetic materials underpin electric motors, generators, transformers, data storage and sensing, and the search for new ones (for example, permanent magnets with less rare earth, or materials for magnetocaloric cooling) is a long-standing priority~\cite{gutfleisch2011magnetic,coey2010magnetism}. High-throughput density functional theory (DFT) databases~\cite{jain2013materials,saal2013oqmd,kirklin2015oqmd,choudhary2020jarvis,curtarolo2012aflow}, text-mined literature compilations~\cite{court2018curie,itani2025nemad} and curated magnetic-structure databases~\cite{gallego2016magndata} now make it possible to train machine-learning (ML) models for magnetic ordering and for the Curie temperature~\Tc{} at scale~\cite{nelson2019curie,belot2023curie,abediorang2026curie}, complementing high-throughput DFT predictions of magnetic order~\cite{horton2019ordering,frey2020oxides}.

Such models are usually reported through one number: a test-set error or score on a random or grouped split. Three further properties matter to anyone who wants to use the predictions to choose which material to make next. First, whether descriptors grounded in physics (for example DFT formation energies, band gaps, volumes, elastic and dielectric tensors) add information beyond composition, as the phrase ``physics-informed'' implies. Second, whether the model knows when it is wrong: whether its prediction intervals and class probabilities are calibrated, overall and for the hard cases. Third, whether accuracy survives a move to chemistry the model has not seen, which is the actual use case in discovery~\cite{meredig2018extrapolation,li2023critical}. Shared benchmarks have standardized how accuracy is compared~\cite{dunn2020matbench}, and work on uncertainty quantification, data leakage and out-of-distribution testing shows that accuracy alone can mislead~\cite{tran2020uq,kapoor2023leakage,li2023critical}. These properties are, however, rarely measured together, on the same data and models, with intervals that account for the dependence between records from the same formula or the same source paper.

Here we measure all three on one curated dataset and one set of models, with comparisons fixed in advance. We ask:
\begin{itemize}
\item \textbf{RQ1.} Do additional physics descriptors improve the prediction of magnetic ordering, \Tc{} and the magnetic moment?
\item \textbf{RQ2.} Are the predictions' uncertainties reliable and calibrated, overall and for subgroups?
\item \textbf{RQ3.} Do the predictions generalize to unfamiliar materials: new compositions, new chemical systems, other label sources?
\end{itemize}
Our central finding is that these are distinct properties. A model that scores well on one can fail on another, so a single accuracy figure is not evidence of physical insight, of calibrated uncertainty, or of usefulness for discovery. We also report the errors we found in our own earlier analyses, because several of them (a string-based novelty test, a default label written by a data collector, a per-atom normalization that used the wrong cell) are likely to recur in other work.

%% ===================================================================================================================
\section{Results}\label{sec:results}

\subsection{A curated multi-source dataset with matched descriptor tiers}\label{sec:dataset}

We merged nine files from seven providers (Materials Project~\cite{jain2013materials}, OQMD~\cite{saal2013oqmd}, JARVIS-DFT~\cite{choudhary2020jarvis}, two NEMAD releases~\cite{itani2025nemad}, the ChemDataExtractor Curie/N\'eel database~\cite{court2018curie}, MAGNDATA~\cite{gallego2016magndata} and two experimental literature compilations) into one schema with row-level provenance (507,828 rows; Methods). Every material received composition, elemental and electronic descriptors. Materials with a DFT counterpart also received 47 harmonized physics properties from Materials Project, JARVIS and AFLOW~\cite{curtarolo2012aflow}.

Predictors were organised into three nested tiers. Tier~1 holds 65 composition-only descriptors, available for every material. Tier~2 adds 8 oxidation-state descriptors (73 in total). Tier~3 adds 24 DFT physics descriptors (97 in total): formation energy and energy above hull, band gap and band edges, Fermi level, space group, crystal system, point group, volume and density, elastic moduli and anisotropy, dielectric constants, refractive index and maximum piezoelectric coefficient. Tier~3 is available for 28,512 materials.

Target curation was the largest single source of change in our results, so we describe it in some detail (Methods). Exclusions followed written rules and only removed rows; no value was ever edited. The text-mined ``Curie temperature'' column contained superconducting transition temperatures, ferroelectric Curie points and N\'eel temperatures. Five rules (no magnetism-capable element; non-ferromagnetic element; $<1$\,K; $>1400$\,K, above the highest known ferromagnetic \Tc{} of Co; formula labelled only antiferromagnetic in every source) reduced the target from 34,322 to 29,740 records. All 15,493 JARVIS ordering labels were the constant ``ferromagnetic'' returned by a collector default, not an assignment, and were quarantined. Table~\ref{tab:table1_dataset} lists the final targets. Three secondary targets (mass and volume magnetization, anisotropy constant $K_1$) are small and are analysed only in the Supplementary Information.

All splits group records by reduced formula, so no composition appears in both training and test data (with the small exception discussed under Limitations). For each target, a \emph{matched} population of materials with complete Tier~3 descriptors is used at every tier: identical rows, splits, hyperparameter grids and metrics. Differences between tiers can then come only from the added columns.

\input{../tables/table1_dataset.tex}

Composition-only (Tier~1) models are already accurate on the full populations. For ordering, XGBoost~\cite{chen2016xgboost} reached macro-F1 0.727 [0.711, 0.744] and histogram gradient boosting~\cite{ke2017lightgbm,pedregosa2011sklearn} 0.720, against 0.469 for multinomial logistic regression and 0.151 for a majority-class reference. For \Tc{}, gradient boosting gave a mean absolute error (MAE) of 76.2\,K [69.9, 82.5] and $R^2$ 0.78, and XGBoost 75.4\,K; Ridge regression gave 147.2\,K and the training-mean reference 218.8\,K. For the moment per cell, the validation-selected gradient-boosting model reached MAE 2.58\,\muB{}/cell ($R^2\approx0.70$). Tree ensembles were statistically tied within every target. Throughout, ``stable'' means that the 95\% bootstrap interval of a difference excludes zero; for regression the bootstrap resamples formula groups, and for the external benchmark it resamples source-paper clusters (Methods).

%% -------------------------------------------------------------------------------------------------------------------
\subsection{RQ1: physics descriptors add little beyond composition}\label{sec:rq1}

On the matched populations we compared Tier~1 with Tier~2 and with Tier~3 over ten pre-specified formula-grouped splits for three algorithms (Table~\ref{tab:table2_physics_descriptors}, Fig.~\ref{fig:rq1}).

\paragraph{Magnetic ordering.} No algorithm gained macro-F1 from either tier (matched $N=6{,}061$). Gradient boosting even lost $-0.029$ [$-0.057$, $-0.001$] from Tier~1 to Tier~2. An apparent Tier~2 improvement of $+0.036$ in our earlier single-split analysis did not replicate once the default ferromagnetic labels were removed and the comparison was repeated over ten splits.

\paragraph{Curie temperature.} No tier gave a practical MAE improvement (matched $N=4{,}602$). For gradient boosting the Tier~1$\rightarrow$3 change was $-14.3$\,K [$-31.3$, $+3.8$]: an improvement larger than 3.8\,K is excluded, while a degradation of up to about 30\,K is not. We therefore report ``no practical improvement (below 10\,K)'', not equivalence. Ridge fits with Tier~3 were numerically unstable. A few non-physical inputs (an elastic-anisotropy ratio of order $10^{110}$, extreme bulk moduli) produced exploded predictions; clipping inputs at training percentiles removes this, so the instability is not evidence about physics. On the full population, Tier~2 produced a small degradation for the tree models (gradient boosting $-3.8$\,K [$-7.1$, $-0.4$]) that was not robust to the antiferromagnet-only exclusion rule.

\paragraph{Magnetic moment.} Here Tier~3 gave a small, reproducible gain per cell (gradient boosting $+0.079$ [$+0.049$, $+0.111$]\,\muB{}/cell; Ridge $+0.055$; XGBoost $+0.026$ to $+0.062$ across thread and seed realisations, 1.3--3.1\% of MAE). The DFT moment is a total per unit cell, so it scales with the number of atoms in the cell. Adding a single scalar, the number of sites in the Materials Project cell, to Tier~2 recovered the gain for all three algorithms in all ten splits (gradient boosting $+0.055$ [$+0.031$, $+0.080$]). What Tier~3 added beyond site count was not stable for any algorithm (gradient boosting $+0.024$ [$-0.003$, $+0.050$]; XGBoost $-0.014$ [$-0.045$, $+0.018$]). For the moment per atom, a modest gain remained: gradient boosting $+2.9\%$ [1.5, 4.4], XGBoost $+2.8\%$ [1.6, 3.9], Ridge $+6.0\%$ [5.1, 6.8]. The gain was larger on subsets where the atom count verifiably belongs to the same calculation as the moment (up to 7.9\% for gradient boosting). Add-one analyses attribute it mainly to outputs of that same DFT calculation, with a smaller share from the volume family. This is a modest secondary result. It replaces an earlier, larger estimate (14\%, 9\% and 6\%), which we withdrew after finding that 25.4\% of rows divided the moment by the atom count of a different cell.

\input{../tables/table2_physics_descriptors.tex}

\begin{figure}[t]
\centering
\includegraphics[width=\textwidth]{../figures/fig1_physics_descriptors.pdf}
\caption{\textbf{Physics descriptors add little beyond composition (RQ1).} Change in test performance from Tier~1 to Tier~2 (oxidation chemistry) and to Tier~3 (DFT physics) on identical matched materials, ten formula-grouped splits, with 95\% intervals; positive values are improvements. \textbf{a}, ordering macro-F1. \textbf{b}, Curie temperature MAE (the unstable Ridge Tier~3 estimate, $-825$\,K, lies off scale and is drawn as an arrow). \textbf{c}, moment per cell. \textbf{d}, moment per atom, relative MAE gain.}\label{fig:rq1}
\end{figure}

\paragraph{Attribution is not benefit.} TreeSHAP~\cite{lundberg2017shap,lundberg2020treeshap} assigns 40\% (ordering) and 36\% (\Tc{}) of the mean absolute attribution of Tier~3 XGBoost models to physics descriptors, although those descriptors do not improve accuracy. When informative composition features are correlated with physics features, the model can route the same information through either set. A large attribution therefore shows that the model uses a feature, not that the feature adds anything. Feature rankings for ordering and the moment were stable over the ten splits (mean pairwise Spearman $\geq0.95$; the magnetic-element fraction ranked first in every split). For \Tc{} the rankings were not stable (Spearman 0.83--0.94, the top feature changed between splits), so we make only family-level statements about \Tc{} attributions.

%% -------------------------------------------------------------------------------------------------------------------
\subsection{RQ2: marginal calibration is attainable; subgroup and shifted coverage is not}\label{sec:rq2}

\paragraph{Regression intervals.} We compared three conformal methods~\cite{vovk2005algorithmic,lei2018distribution,angelopoulos2023gentle} at 90\% nominal coverage. These were constant-width split conformal, normalized conformal with a learned difficulty scale~\cite{papadopoulos2011normalized}, and conformalized quantile regression (CQR)~\cite{romano2019cqr}. Each was calibrated on a held-out 20\% of training formulas (Table~\ref{tab:table3_regression_intervals}, Fig.~\ref{fig:rq2}a,b). For \Tc{}, only normalized conformal stayed close to nominal across splits: 90.3\% mean coverage (SD 2.5), with narrower intervals (366\,K mean width on the frozen split, against 552\,K for split conformal). Split conformal and CQR varied more from split to split (SD 6.0 and 5.8 points). On the hardest decile of test materials, ranked by the model's own difficulty estimate, split conformal covered 56\% and CQR 63\%, against 90\% for normalized conformal. For the moment, all three methods were within 0.1 points of nominal in every split. Split conformal still covered only half of the hardest decile (49.7\%) and 84\% of oxides. CQR failed where 38\% of Materials Project rows have exactly zero moment: calibrated on JARVIS and tested on Materials Project, it covered 35\%, because a negative quantile correction excluded the zeros.

Coverage is also not protected under a change of label source. When intervals for \Tc{} were calibrated on curated rows and tested on text-mined rows, coverage fell to 76.9\% (split), 80.3\% (CQR) and 85.1\% (normalized).

\input{../tables/table3_regression_intervals.tex}

\paragraph{Ordering probabilities and prediction sets.} The ordering classifier was trained with balanced class weights, so its raw probabilities reflect a balanced prior. Their expected calibration error (ECE) was 0.059. Correcting them to the true class prior~\cite{saerens2002adjusting} made calibration worse (0.093) until a temperature was also fitted ($T=1.65$; ECE 0.016)~\cite{guo2017calibration} (Table~\ref{tab:table4_ordering_calibration}). Marginal conformal prediction sets met 90\% coverage overall (90.0\% over ten splits) by covering the majority classes, which hid the failure on the minority classes. They covered ferromagnets 95.6\% and non-magnets 97.3\% of the time, but antiferromagnets only 57.7\% and ferrimagnets 52.7\%. Class-conditional (Mondrian) sets~\cite{vovk2013conditional} covered 89.9--91.1\% for every class over ten splits, at the cost of a larger mean set size (1.57 against 1.34 classes). Sets calibrated on Materials Project labels and tested on experimental labels covered 83.5\%.

\input{../tables/table4_ordering_calibration.tex}

\begin{figure}[t]
\centering
\includegraphics[width=\textwidth]{../figures/fig2_uncertainty_calibration.pdf}
\caption{\textbf{Marginal calibration hides subgroup failures (RQ2).} \textbf{a,b}, 90\% interval coverage over ten splits, all test materials against the hardest decile, for the Curie temperature (\textbf{a}) and the moment per cell (\textbf{b}) under split, normalized and quantile (CQR) conformal prediction. \textbf{c}, per-class coverage of 90\% ordering prediction sets, marginal against class-conditional calibration. The dashed line marks nominal coverage.}\label{fig:rq2}
\end{figure}

%% -------------------------------------------------------------------------------------------------------------------
\subsection{RQ3: generalization weakens with chemical distance}\label{sec:rq3}

\paragraph{Internal tests.} Formula-grouped splits still let a test material share its chemical system (its set of elements) with training materials; this was the case for 65--78\% of test rows. When whole chemical systems were held out instead, ordering macro-F1 fell from 0.712 to 0.684 ($-0.028$ [$-0.041$, $-0.016$]; worse in all ten paired splits), with ferrimagnet recall falling most (Table~\ref{tab:table6_generalization_internal}). The moment degraded slightly but not stably (+0.073\,\muB{}/cell [$-0.009$, $+0.154$]), and the \Tc{} test was too noisy to conclude ($+2.7$\,K [$-16.8$, $+22.2$]).

Ordering skill did not transfer across label provenance. A model trained on Materials Project labels and tested on experimental labels scored macro-F1 0.284, with balanced accuracy 0.323, below the three-class chance level of 0.333. The reverse direction scored 0.159. Within each provenance the same algorithm scored 0.637 (computed) and 0.784 (experimental). The pooled headline score therefore does not describe a general ability to predict magnetic order, but agreement with the labelling convention of each source.

\input{../tables/table6_generalization_internal.tex}

\paragraph{External benchmark.} We scored the Tier~1 XGBoost \Tc{} model and its normalized conformal intervals, without retraining or tuning, on an independent public table of 2,500 experimental Curie temperatures with source DOIs, curated for a separate ML study of \Tc{}~\cite{malaei_magnetic_data,abediorang2026curie}. Before interpreting any error, we audited the table under a written plan (Methods; Supplementary Table~S1). The audit changed the picture substantially. By composition, not by formula string, 1,998 of the 2,500 records (80\%) were already in our repository, and 1,200 (48\%) carried a value identical to within 0.5\,K to one we held for the same composition. Both datasets draw on the same upstream compilations. On those records the model scores $R^2$ 0.90, which reflects recall of training labels, not prediction. We excluded 105 records for reasons fixed before looking at errors: 60 computational labels, 18 malformed compositions, and 35 values outside our target definition (mostly N\'eel temperatures of antiferromagnets). This left a primary benchmark of 2,395 records.

On the 465 compositions that the model had truly never seen, the MAE was 102.7\,K [78.7, 154.3], $R^2$ 0.69, and the 90\% intervals covered 89.5\% (Table~\ref{tab:table5_external_benchmark}, Fig.~\ref{fig:rq3}). Intervals come from a bootstrap over source-paper clusters, because records from one paper are not independent. Blind compositions in chemical systems present in training had MAE 90.6\,K and $R^2$ 0.78. The 69 blind records in unseen chemical systems had MAE 171.9\,K and $R^2$ $-0.17$ [$-0.93$, 0.66], with $R^2<0$ in 70\% of resamples. These 69 records come from only 26 papers. Before the analysis we had fixed a wording rule: ``no skill beyond the mean'' would be claimed only if the upper $R^2$ limit was at most 0.2. It is not, so the supported statement is that accuracy in unseen chemical systems is much weaker, with wide uncertainty. The drop in $R^2$ relative to seen systems is resolved under every clustering ($-0.95$ [$-1.69$, $-0.06$] with paper clusters). The MAE gap ($+81$\,K) is resolved when resampling chemical systems but not source papers (Supplementary Table~S2).

\input{../tables/table5_external_benchmark.tex}

\begin{figure}[t]
\centering
\includegraphics[width=\textwidth]{../figures/fig3_external_benchmark.pdf}
\caption{\textbf{External benchmark of experimental Curie temperatures (RQ3).} MAE (\textbf{a}) and $R^2$ (\textbf{b}) for nested novelty subsets of the external table, before and after the audit, with 95\% intervals from a bootstrap over source-paper clusters. ``Blind'' means the composition never occurs in the training repository; ``unseen system'' means its set of elements never occurs either.}\label{fig:rq3}
\end{figure}

\paragraph{Chemical distance.} Before comparing it with any error, we fixed a similarity measure: the total-variation distance between element-fraction vectors, to the nearest training formula. On the blind records, error rose with this distance (Spearman $\rho=0.28$ [0.10, 0.49]); MAE increased from 62\,K in the closest quartile to 165\,K in the farthest ($+103$\,K [$+31$, $+176$]). Normalized conformal intervals widened with distance (388 to 529\,K mean width), but not enough. Coverage fell from 0.95 to 0.85 ($-0.10$ [$-0.23$, $-0.01$]; Fig.~\ref{fig:distance}, Supplementary Table~S3). Distance explains only part of the failure in unseen chemical systems. Blind records in seen systems at comparable distance still scored MAE 109\,K and $R^2$ 0.75, against 172\,K and $-0.17$ in unseen systems ($R^2$ difference $-0.92$ [$-1.66$, $-0.05$]).

\begin{figure}[t]
\centering
\includegraphics[width=\textwidth]{../figures/fig4_error_vs_distance.pdf}
\caption{\textbf{Error grows, and coverage falls, with distance from the training chemistry (RQ3).} For the 465 blind external records, binned by composition distance to the nearest training formula (quartiles): MAE (\textbf{a}), mean width of the 90\% interval (\textbf{b}) and its coverage (\textbf{c}), with 95\% paper-cluster bootstrap intervals.}\label{fig:distance}
\end{figure}

%% ===================================================================================================================
\section{Discussion}\label{sec:discussion}

The four properties we measured came apart in every pairing we could test. \emph{Accuracy versus physics enrichment}: composition alone was already accurate, DFT descriptors gave no practical gain for ordering or \Tc{}, and they still received more than a third of the SHAP attribution. \emph{Accuracy versus calibration}: the most accurate ordering model had raw probabilities that became miscalibrated under the correct class prior, and its marginal 90\% sets covered antiferromagnets and ferrimagnets little more than half the time. \emph{Calibration versus generalization}: intervals that were calibrated on blind external compositions lost coverage with chemical distance and under a change of label source. \emph{Accuracy versus generalization}: a \Tc{} model with $R^2$ 0.78 within the training distribution had a negative point estimate of $R^2$ in unseen chemical systems. An ordering model with macro-F1 0.73 was at chance across label provenance. None of these properties can be inferred from another, so each has to be measured directly.

\paragraph{Why physics descriptors did not help.} Several explanations are consistent with our results, and we cannot separate them. Composition fixes much of the information that the DFT descriptors encode, so tree ensembles may already recover it from Tier~1. The DFT descriptors describe the lowest-energy computed polymorph of a formula, while the targets, especially experimental \Tc{}, may refer to a different structure, a doped sample or a thin film. And the matched populations are restricted to materials with a full DFT record, a selection that is not random. Where a gain did appear (the moment), it traced back to the size of the cell and to outputs of the same calculation that produced the target, not to independent physics. We therefore caution against reading SHAP attributions on physics features as evidence that physics was learned. We also caution that ``physics-informed'' gains need matched rows, repeated splits and a simple control (here, the number of sites in the cell) before they are claimed.

\paragraph{Practical recommendations.} Our results suggest six practices for magnetic-property ML, and likely for other property-prediction tasks:
(i) compare descriptor sets on identical rows and repeated grouped splits, and report intervals rather than single-split differences;
(ii) report normalized or class-conditional conformal coverage by subgroup (hardest decile, minority class, family, label source), not only marginally;
(iii) define novelty by composition identity and chemical system, not by formula strings, and check whether external ``test'' values duplicate training labels;
(iv) resample at the level of the dependence structure (formula, source paper, chemical system) when computing intervals;
(v) audit data collectors for default values (here, a constant ferromagnetic label for every JARVIS entry); and
(vi) state claims across seeds when results depend on thread count. For XGBoost, the thread count acted like a hidden random seed: the mean prediction changed by about 25\% of the MAE between thread settings.

\paragraph{Limitations.} The external benchmark is a text-mined compilation that overlaps heavily with our sources. After correction only 465 records are truly blind, and the conclusions about unseen chemical systems rest on 69 records from 26 papers, which is why their intervals are wide. Formula-grouped splits are not perfectly disjoint: some compositions are spelled in two ways (fractional and integer formulas), which affects 0.6\% of \Tc{} test rows and 1.3\% of ordering test rows (none for the moment). This is too small to change our conclusions. Our target-curation rules are screens with known false positives. For example, Ti$^{3+}$ ferromagnets such as YTiO$_3$ are removed by the magnetism-capable-element rule, and the ferroelectric/N\'eel ambiguity of multiferroics cannot be detected from composition alone. We studied gradient-boosted trees and linear models on tabular descriptors only; graph neural networks on crystal structures might use physics information differently. Our conclusions about Tier~3 apply to the descriptors we tested, not to every possible physics descriptor. Finally, the redistribution licences of the upstream data have not all been verified (see Data availability).

\paragraph{Conclusion.} On a curated dataset of more than 100,000 labelled magnetic materials, composition-only models predicted ordering, \Tc{} and the magnetic moment well within the training distribution. DFT physics descriptors added little once comparisons were matched and controlled. Calibration was attainable overall but not for the hardest materials, the minority magnetic classes or shifted label sources. Accuracy fell substantially in unfamiliar chemistry. For materials discovery, where the candidates of interest are by definition unfamiliar, a model's in-distribution accuracy is the least informative of the four properties, and calibrated distance-aware uncertainty is the most needed.

%% ===================================================================================================================
\section{Methods}\label{sec:methods}

\subsection{Data collection and harmonization}
Sources were Materials Project~\cite{jain2013materials} (100,000 entries with computed ordering and total magnetization), OQMD~\cite{saal2013oqmd,kirklin2015oqmd} via jarvis-tools, JARVIS-DFT~\cite{choudhary2020jarvis} (entries with magnetic moment $>0.05$\,\muB{}), two NEMAD releases~\cite{itani2025nemad}, the ChemDataExtractor Curie/N\'eel database~\cite{court2018curie}, MAGNDATA~\cite{gallego2016magndata}, and two experimental literature compilations of magnetic and anisotropy properties. All sources were mapped to a common 25-column schema. Each numeric property was stored with its value, unit and measurement temperature, and every row carried its source file. Records were deduplicated on material name plus DOI, not on formula alone, so that independent measurements of one composition were kept (507,828 rows). Formulas were canonicalized with pymatgen~\cite{ong2013pymatgen} to reduced formulas, and 1,940 records with non-chemical ``formulas'' (mineral or trade names, generic templates such as ABO$_3$) were quarantined.

\subsection{Descriptors}
Tier~1 (65 columns) contains composition statistics in the spirit of Magpie~\cite{ward2016magpie}: element-presence flags and fractions for the main magnetic and light elements; transition-metal, rare-earth, actinide and magnetic-element fractions; $s/p/d/f$-block fractions; average unpaired-electron, $d$- and $f$-electron counts (counted outside the noble-gas core); valence-electron concentration; means and variances of atomic radius, electronegativity, mass, melting temperature, period and group; size mismatch; and mixing and formula entropies. Tier~2 adds 8 oxidation-state statistics from pymatgen's charge-balancing guesser, including a flag for whether a solution was found. Tier~3 adds 24 DFT properties joined by reduced formula, taking the most stable polymorph. Materials Project was the primary source for most properties; JARVIS was primary for elastic and dielectric tensors where it is denser, and AFLOW~\cite{curtarolo2012aflow} for thermal properties. Crystal system and point group were treated as categorical variables. Finite values that overflow single precision (four rows) were set to missing at load time.

\subsection{Target definitions and curation}
\emph{Magnetic ordering.} 116 raw labels were mapped to four classes (FM, AFM, FiM, NM). Labels with qualifiers that do not establish a class (for example ``ferromagnetic-like'', ``nearly ferromagnetic'') and distinct states (superparamagnetic, spin glass) were excluded; compensated ferrimagnets were kept as FiM. All JARVIS rows were quarantined because their label was a constant written by the collector (rule O1). \emph{Curie temperature.} Values were standardized to kelvin. Zeros that a single source used as placeholders were removed, as were 54 encoded values. Rules C1--C5 then removed formulas without a magnetism-capable element (C1), single-element formulas that are not ferromagnetic elements (C2), values below 1\,K (C3) or above 1400\,K (C4), and formulas labelled only AFM in every source (C5). A sensitivity target without C5 ($N=30{,}059$) gave the same baselines. \emph{Magnetic moment.} The Materials Project total magnetization is a DFT total moment per unit cell at 0\,K. Moments in mass or volume units were kept as separate secondary targets, never converted into one another, and implausible values were removed (rule set v1.4). The per-atom target divides the moment by the number of sites of the cell of the same calculation. A first version used site counts from a different cell for 25.4\% of rows; this was corrected before the results reported here. All exclusions are listed with reason codes in the data package.

\subsection{Splits}
Records were assigned to training, validation and test sets (70/15/15) by reduced-formula group, using a size-balanced, stratified group assignment (strata: classes, or quantile bins of the formula-mean target). We did not use StratifiedGroupKFold, because a few formulas with many records (for example Fe, Mn, Co, Ni) made its partitions badly unbalanced. Robustness analyses used ten further splits with pre-specified seeds ($20267000+i$). Leave-chemical-system-out splits grouped records by their element set instead of their formula.

\subsection{Models and selection}
We used multinomial logistic regression or Ridge regression (median imputation, standardization), scikit-learn histogram gradient boosting~\cite{pedregosa2011sklearn,ke2017lightgbm} and XGBoost~\cite{chen2016xgboost} (histogram trees). Both boosters handle missing values natively. Classifiers used balanced class weights; regression targets were used in raw units. Hyperparameters were chosen from small fixed grids by validation macro-F1 (classification) or RMSE (regression), refitted on training plus validation data, and scored on the test set once. The grids were: $C\in\{0.1,1,10\}$ or $\alpha\in\{0.1,1,10,100\}$ for the linear models; 500 iterations with learning rate $\in\{0.05,0.1\}$ and maximum leaf nodes $\in\{31,127\}$ for gradient boosting; and 500 trees, depth $\in\{6,10\}$, learning rate $\in\{0.05,0.1\}$ and subsample and column fractions of 0.8 for XGBoost. Reference models predicted the majority class or the training mean. XGBoost results depend on the number of threads, so all XGBoost claims are stated across splits or seeds.

\subsection{Statistical comparison}
Tier differences were computed on identical test materials, with the same formula-group bootstrap draws for both tiers (paired; 1,000 resamples). They are reported per split, as a mean over the ten splits with a formula-bootstrap interval, and as across-split sign counts. A difference was called stable when its 95\% interval excluded zero. For \Tc{}, practical relevance used a 10\,K margin. Comparisons, populations, metrics and decision thresholds were written into configuration files before the corresponding models were fitted, and later deviations were logged.

\subsection{Uncertainty quantification}
Whole formulas (20\% of the training formulas) were set aside as a calibration set; models were refitted on the remaining 80\%. \emph{Split conformal} used absolute residuals. \emph{Normalized conformal} divided residuals by a difficulty estimate from a cross-fitted model of the absolute residual. \emph{CQR} used XGBoost quantile models at 5\% and 95\%. For ordering, we compared raw probabilities, prior-corrected probabilities~\cite{saerens2002adjusting}, and prior-corrected probabilities with a temperature fitted on the calibration set~\cite{guo2017calibration}. ECE used 15 equal-width confidence bins. Prediction sets were built with a marginal or a class-conditional threshold on one minus the probability of the true class~\cite{vovk2013conditional}. Coverage was reported overall, per difficulty decile, per material family, per label source and per class.

\subsection{External benchmark and audit}
The external table~\cite{malaei_magnetic_data,abediorang2026curie} lists 2,500 ferromagnets with experimental \Tc{} and source DOIs. It is not redistributed; a script downloads the analysed version and checks its SHA-256 hash. The audit plan (checks A1--A8, Supplementary Table~S1) was written before any error was examined, and three amendments were logged. Novelty was defined by composition identity against the training repository, never by formula string, and source papers were matched on normalized DOIs. Composition distance is the total-variation distance between element-fraction vectors (0 for identical compositions, 1 for no shared element), measured to the nearest of the 23,757 training formulas. A z-scored Tier~1 descriptor distance was used as a secondary measure ($\rho=0.21$ [0.01, 0.31]). Intervals come from a cluster bootstrap (10,000 resamples) over source-paper clusters, defined as records connected by any shared DOI. As sensitivity analyses we used first-listed-DOI clusters and chemical-system clusters. Removing the 124 records that come only from papers on films, nanoparticles, high pressure or 2D materials changed subset MAEs by at most 3.3\,K.

\subsection{Errors found and corrected during the study}
We report these because they changed headline numbers and are easy to repeat. (1) The collector default ``ferromagnetic'' label for JARVIS entries (16\% of an earlier ordering target). (2) Non-Curie transition temperatures in a text-mined \Tc{} column. (3) Per-atom moments divided by the atom count of a different cell (25.4\% of rows). (4) An external ``blind'' test that compared formula strings, so it counted compositions we already held as new (1,302 claimed blind records against 502 by composition). (5) A DOI normalizer that missed the \url{https://doi.org/} prefix. (6) A single-split Tier~2 improvement for ordering that did not replicate. Earlier results files are kept unchanged, and each correction is recorded in a versioned claims register.

%% ===================================================================================================================
\backmatter

\bmhead{Supplementary information}
Supplementary Tables S1--S3 (external-benchmark audit, subset comparisons under alternative clusterings, error by composition distance), Supplementary Fig.~S1 (absolute error against distance for every blind record), secondary-target analyses (mass and volume magnetization, $K_1$), material-family analyses, bias-reduction experiments and the versioned claims register. \todo{assemble the Supplementary Information document from ../tables/tableS*.tex, ../figures/figS1 and ../results/supplementary/}

\bmhead{Acknowledgements}
\todo{funding and acknowledgements}

\section*{Declarations}

\begin{itemize}
\item \textbf{Funding:} \todo{}
\item \textbf{Competing interests:} The authors declare no competing interests. \todo{confirm}
\item \textbf{Data availability:} The curated analysis dataset (versions 1.2--1.4), the frozen splits, the external-benchmark adjudications and all result files are provided in the reproducibility package \todo{repository and DOI, e.g. Zenodo}. The dataset is derived from Materials Project, OQMD, JARVIS-DFT, AFLOW, NEMAD, the ChemDataExtractor Curie/N\'eel database, MAGNDATA and two experimental compilations, each under its own terms. \todo{redistribution licences of every upstream source must be verified before deposit}. The external benchmark table is not redistributed because no licence is stated; a script retrieves the analysed version and verifies its hash. \todo{three result files contain its compound names and values: clear the licence or drop those columns}
\item \textbf{Code availability:} All analysis code, configuration files with fixed seeds, table and figure builders, and tests are included in the reproducibility package \todo{repository URL and licence}.
\item \textbf{Author contributions:} \todo{CRediT statement, e.g. M.I.: conceptualization, data curation, methodology, software, formal analysis, writing (original draft); G.P.H.: supervision, methodology, writing (review and editing). Confirm with both authors.}
\item \textbf{Use of AI tools:} \todo{The data pipeline, analyses and this draft were prepared with substantial assistance from an AI coding assistant (Claude, Anthropic). State the venue-required disclosure; the authors take responsibility for all content. Draft with /ars-disclosure for the target journal.}
\end{itemize}

\bibliography{references}

\end{document}
